% Part: first-order-logic
% Chapter: introduction
% Section: soundness-completeness

\documentclass[../../../include/open-logic-section]{subfiles}

\begin{document}

\olfileid{fol}{int}{scp}

\olsection{યથાર્થતા અને પૂર્ણતા}

પ્રથમ-ક્રમ તર્કશાસ્ત્ર માટે !!{derivation} તંત્રો પણ રજૂ કરીશું.
તર્કશાસ્ત્રીઓએ ઘણાં !!{derivation} તંત્રો વિકસાવ્યાં છે, પરંતુ એ બધાં
!!{sentence}s વચ્ચેનો એક જ !!{derivability} સંબંધ વ્યાખ્યાયિત કરે છે.
ચોક્કસ રીતે વ્યાખ્યાયિત કોઈ પ્રકારની !!a{derivation} હોય, તો આપણે
કહીએ છીએ કે~$\Gamma$,~$!A$ને \emph{!!{derive}s} કરે છે અને
$\Gamma \Proves !A$ લખીએ છીએ. !!^{derivation}s હંમેશાં સંકેતોની
સાન્ત ગોઠવણો હોય છે---કદાચ !!{sentence}sની યાદી અથવા કોઈ વધુ જટિલ
રચના. !!{derivation} તંત્રોનો હેતુ કોઈ ગણ~$\Gamma$થી !!a{sentence}
ફલિત થાય છે કે નહિ તે નક્કી કરવાનું સાધન આપવાનો છે. એ હેતુ પાર પાડવા
$\Gamma \Entails !A$ ત્યારે, અને માત્ર ત્યારે જ, $\Gamma \Proves !A$
હોવું જોઈએ.

જો $\Gamma \Proves !A$ હોય પરંતુ $\Gamma \Entails !A$ ન હોય, તો
આપણું !!{derivation} તંત્ર અતિશય પ્રબળ હશે અને વધારે પડતું સાબિત કરશે.
$\Gamma \Proves !A$ હોય તો $\Gamma \Entails !A$ હોય એ ગુણધર્મને
\emph{યથાર્થતા} કહે છે; કોઈ પણ સારા !!{derivation} તંત્ર માટે આ લઘુતમ
શરત છે. બીજી તરફ, જો $\Gamma \Entails !A$ હોય પરંતુ $\Gamma \Proves
!A$ ન હોય, તો આપણું !!{derivation} તંત્ર અતિશય નબળું છે અને પૂરતું
સાબિત કરતું નથી. $\Gamma \Entails !A$ હોય તો $\Gamma \Proves !A$
હોય એ ગુણધર્મને \emph{પૂર્ણતા} કહે છે. યથાર્થતા સામાન્ય રીતે પ્રમાણમાં
સહેલાઈથી સાબિત થાય છે (અનુમાનાત્મક રીતે વ્યાખ્યાયિત
!!{derivation}sની રચના પર અનુમાન વડે). પૂર્ણતા સાબિત કરવી વધુ મુશ્કેલ
છે.

યથાર્થતા અને પૂર્ણતાનાં અનેક મહત્ત્વનાં પરિણામો છે. !!{sentence}sનો
ગણ~$\Gamma$ કોઈ વ્યાઘાત (જેમ કે $!A \land \lnot !A$) !!{derive}s કરે,
તો તેને \emph{અસુસંગત} કહે છે. અસુસંગત~$\Gamma$ને કોઈ નિદર્શ હોઈ
શકતો નથી; તે અસંતોષ્ય છે. પૂર્ણતામાંથી ઊલટું વિધાન ફલિત થાય છે: જે
કોઈ~$\Gamma$ અસુસંગત નથી---અથવા, આપણે કહીએ તેમ,
\emph{સુસંગત} છે---તેને નિદર્શ છે. હકીકતમાં, આ વાત પૂર્ણતાને સમતુલ્ય
છે અને આપણે ખરેખર પૂર્ણતાનું આ જ સ્વરૂપ સાબિત કરીશું. આ ઊંડું અને
કદાચ આશ્ચર્યજનક પરિણામ છે:~$\Gamma$માંથી $!A \land \lnot !A$ સાબિત
કરી શકાતું નથી એટલી જ વાત ખાતરી આપે છે કે~$\Gamma$ જેવું વર્ણન કરે
છે તેવી કોઈ !!a{structure} અસ્તિત્વમાં છે. તેથી પૂર્ણતા આ પ્રશ્નનો
જવાબ આપે છે: !!{sentence}sના કયા ગણોને નિદર્શો છે? જવાબ: બરાબર બધા
સુસંગત ગણોને.

યથાર્થતા અને પૂર્ણતા પ્રમેયોનાં બે મહત્ત્વનાં પરિણામો છે: સઘનતા અને
લેવેનહાઇમ--સ્કોલેમ પ્રમેયો. નિદર્શોના સિદ્ધાંતમાં આ મહત્ત્વનાં
પરિણામો છે અને બીજાં ઘણાં રસપ્રદ પરિણામો સ્થાપવા વાપરી શકાય છે. એમાંનાં
બેનો ઉલ્લેખ આપણે કરી ચૂક્યા છીએ: !!a{structure}નું !!{domain} સાન્ત છે
અથવા !!{nonenumerable} છે એ પ્રથમ-ક્રમ તર્કશાસ્ત્ર વ્યક્ત કરી શકતું
નથી.

ઐતિહાસિક રીતે આ બધું---પ્રથમ-ક્રમ તર્કશાસ્ત્રની વાક્યરચના અને
અર્થવિચાર કેવી રીતે વ્યાખ્યાયિત કરવાં, સારાં !!{derivation} તંત્રો
કેવી રીતે વ્યાખ્યાયિત કરવાં, તેઓ યથાર્થ અને પૂર્ણ છે એવું કેવી રીતે
સાબિત કરવું, પ્રથમ-ક્રમ ભાષાઓમાં શું વ્યક્ત કરી શકાય અને શું નહિ એ
સ્પષ્ટ કરવું---સમજવામાં અને બરાબર ઘડવામાં ઘણો સમય લાગ્યો હતો. હવે આ
બધું કેવી રીતે કરવું તે આપણે જાણીએ છીએ, છતાં બધી વિગતોમાંથી પસાર થવું
હજી ગૂંચવણભર્યું અને કંટાળાજનક થઈ શકે. પરંતુ એ મહત્ત્વનું પણ છે,
કારણ કે પ્રથમ-ક્રમ તર્કશાસ્ત્રની રૂપલક્ષી ભાષા માટે અહીં વિકસાવેલી
પદ્ધતિઓનો તર્કશાસ્ત્ર, કમ્પ્યુટર વિજ્ઞાન અને ભાષાશાસ્ત્રમાં વ્યાપક
ઉપયોગ થાય છે. એટલે લાંબા ગાળે વિગતોનો ઝીણવટથી અભ્યાસ ફળ આપે છે.

\end{document}
